\chapter{Angioplasty and Stenting}

\section*{Patient focus}
This chapter explains the examination or clinical problem in accessible language. Patients should use it to prepare questions and understand the purpose, experience, benefits, and limitations of imaging; local instructions and the official report take priority.


\label{ch:ch79_angioplasty}

Angioplasty and stenting are endovascular procedures that open blocked or narrowed blood vessels using balloons and metal stents. They are performed under fluoroscopic (X-ray) guidance and are alternatives to open vascular surgery for many patients.

---

\section{1. Patient Perspective: What Angioplasty Involves}

If one of your blood vessels is narrowed or blocked, you may be offered a procedure to open it up.

\begin{itemize}
  \item A thin tube (catheter) is inserted through a small puncture in the groin (femoral artery) or wrist (radial artery).
  \item Under X-ray guidance, the catheter is guided to the blocked vessel. Contrast dye is injected to visualise the blockage.
  \item A small balloon is inflated at the site of the narrowing, stretching the vessel open.
  \item In many cases, a small mesh tube (stent) is placed to keep the vessel open long-term.
  \item The procedure is usually performed under local anaesthetic with mild sedation. Most patients go home the same day or the next day.
  \item You will need to take blood-thinning medication (aspirin and/or clopidogrel) for several months afterwards to prevent the stent from clotting.
\end{itemize}

---

\section{2. Indications by Vessel}

\begin{table}[h]
\centering
\begin{tabular}{l l l}
\hline
\textbf{Vessel Territory} & \textbf{Indication} & \textbf{Key Consideration} \\
\hline
Carotid artery & Symptomatic stenosis $>$ 50\% (or $>$ 70\% asymptomatic) & Compete with carotid endarterectomy (CEA) \\
Coronary artery & ACS, stable angina refractory to medical therapy & Cardiology-led; covered in cardiac catheterisation chapters \\
Iliac/femoral artery & Claudication, critical limb ischaemia & TASC A/B lesions favour endovascular; TASC C/D may favour surgery \\
Renal artery & Atherosclerotic or fibromuscular dysplasia causing HTN or renal decline & Medical therapy first; revascularise for flash pulmonary oedema or refractory HTN \\
Subclavian artery & Subclavian steal syndrome, arm claudication & Usually stenting preferred over plain angioplasty \\
\hline
\end{tabular}
\caption{Angioplasty Indications by Vessel Territory}
\end{table}

---

\section{3. Pre-Procedure Planning}

\begin{itemize}
  \item \textbf{Cross-sectional imaging}: Review CTA or MRA to assess lesion length, calcification, vessel diameter, and runoff (for peripheral disease).
  \item \textbf{Access}: Femoral (common femoral artery), radial (for carotid/subclavian), or brachial approach.
  \item \textbf{Lesion characteristics}: Short focal stenosis ($<$ 10 cm for peripheral) responds best to angioplasty. Long, heavily calcified, or chronic total occlusions (CTOs) are more complex.
  \item \textbf{Renal function}: Baseline eGFR; consider volume expansion and minimising contrast.
  \item \textbf{Antiplatelet therapy}: Aspirin loading dose pre-procedure; clopidogrel loading if stenting planned.
  \item \textbf{Contrast}: Iodinated contrast; estimate volume needed. CO$_2$ angiography may be used in severe renal impairment.
\end{itemize}

\section{4. Angiographic Findings and Technique}

\begin{itemize}
  \item \textbf{Diagnostic angiography}: The baseline study defines the stenosis length, severity, and runoff.
  \item \textbf{Pre-dilation}: A balloon catheter is inflated across the stenosis. Balloon size is chosen to match the reference vessel diameter.
  \item \textbf{Stenting}: If residual stenosis $>$ 30\% or flow-limiting dissection occurs after balloon angioplasty, a stent is deployed.
  \item \textbf{Stent types}:
    \begin{itemize}
      \item \textbf{Bare-metal stents} (BMS): Standard for peripheral and renal arteries.
      \item \textbf{Drug-eluting stents} (DES): Reduce restenosis; used in coronary and increasingly in peripheral arteries.
      \item \textbf{Covered stents} (stent-grafts): For aneurysmal disease or vessel perforation.
    \end{itemize}
  \item \textbf{Post-dilation}: Balloon inflation within the stent to ensure full expansion.
  \item \textbf{Completion angiography}: Confirms technical success and assesses runoff.
\end{itemize}

\begin{pearlbox}[Balloon Types]
\textbf{Non-compliant balloons} are used for post-dilation within stents (high pressure, maintain shape). \textbf{Semi-compliant balloons} are used for pre-dilation (more flexible, conform to vessel). \textbf{Drug-coated balloons} (DCB) deliver antiproliferative agents (paclitaxel) to reduce restenosis in peripheral arteries.
\end{pearlbox}

---

\section{5. Complications and Management}

\begin{table}[h]
\centering
\begin{tabular}{l l}
\hline
\textbf{Complication} & \textbf{Management} \\
\hline
Access-site haematoma & Manual compression, ultrasound-guided repair, vascular surgery if needed \\
Acute thrombosis & Aspiration thrombectomy, catheter-directed thrombolysis \\
Contrast nephropathy & Volume expansion, avoid nephrotoxins \\
Distal embolisation & Thrombectomy (aspiration or mechanical), catheter-directed thrombolysis \\
Flow-limiting dissection & Extend stent coverage \\
Vessel perforation & Covered stent, prolonged balloon inflation \\
\hline
\end{tabular}
\caption{Complications of Angioplasty and Stenting}
\end{table}

\section{6. Restenosis Surveillance}

\begin{itemize}
  \item \textbf{Restenosis} is the most common long-term complication, occurring in 20--40\% of peripheral interventions within 12 months.
  \item \textbf{Surveillance protocol}: Duplex ultrasound at 1, 6, and 12 months post-procedure, then annually.
  \item \textbf{Restenosis criteria} on duplex: Peak systolic velocity ratio $>$ 2.0 (suggests $>$ 50\% in-stent restenosis).
  \item \textbf{Repeat intervention}: Balloon angioplasty within the stent, drug-coated balloon, or cutting balloon for in-stent restenosis.
  \item \textbf{Risk factor modification}: Smoking cessation, statin therapy, glycaemic control, and antiplatelet adherence are critical.
\end{itemize}

\begin{pearlbox}[Clinical Pearl]
A patient with worsening claudication or rest pain after peripheral stenting should be investigated urgently with duplex ultrasound and potentially CTA. Acute in-stent thrombosis is a vascular emergency.
\end{pearlbox}

---

\section{7. Structured Reporting Example}

\begin{testbox}[Indications:]
68-year-old male with lifestyle-limiting claudication (walking distance 100 metres). ABI 0.6 on the right. CTA shows a 7 cm stenosis of the right superficial femoral artery (SFA).
\end{testbox}
\begin{refbox}[Findings:]
\textbf{Pre-procedure angiography}: A 7 cm stenosis of the mid-right SFA with approximately 85\% diameter reduction. Heavy calcification. Collateral filling via the deep femoral artery. The popliteal artery and three-vessel runoff to the foot are patent.\\
\textbf{Procedure}: Right common femoral artery access (6 Fr sheath). A 6 mm $\times$ 15 cm drug-coated balloon (DCB) was inflated across the lesion at 8 atm for 3 minutes. Post-dilation angiography shows a residual stenosis of $<$ 10\%. No dissection, no thrombus, no distal embolisation. Good pulsatile flow to the foot.\\
\textbf{Post-procedure}: Sheath removed, manual compression applied. ABI post-procedure: 0.95 (right).
\end{refbox}
\begin{pearlbox}[Impression:]
1. Successful angioplasty of a 7 cm right SFA stenosis using a drug-coated balloon.
2. No complications. ABI normalised post-procedure.
3. Surveillance duplex ultrasound at 1 month recommended. Continue dual antiplatelet therapy.
\end{pearlbox}

\clearpage
